\chapter{1930:Hans Fischer}

\label{chap:chemistry1930}

The Nobel Prize in Chemistry for the year 1930 was awarded to Hans Fischer.

\textbf{Motivation:}

for his researches into the constitution of haemin and chlorophyll and especially for his synthesis of haemin

\section{Hans Fischer}

\subsection*{Early Life and Education}

Hans Fischer was born on 1881-07-27 in Höchst am Main, Germany.
He was a German organic chemist awarded the Nobel Prize in Chemistry in 1930 for his researches into the constitution of haemin and chlorophyll and especially for his synthesis of haemin.

\subsection*{Career and Major Research}

Fischer studied chemistry and medicine in Lausanne, Munich, and Marburg, and received his doctorate at the University of Marburg in 1904 under Theodor Zincke.
He held professorships of medical chemistry at the University of Innsbruck from 1916 to 1918 and at the University of Vienna from 1918 to 1921.
In 1921 he succeeded Heinrich Wieland as professor of organic chemistry at the Technische Hochschule in Munich, where he remained for the rest of his life.
His research concerned the pigments of blood, bile, and plants, all of which are built from pyrrole rings.
He determined the constitution of the bile pigments bilirubin and biliverdin and of haemin, the iron-containing pigment of haemoglobin, and he worked extensively on chlorophyll.
In 1929, together with Karl Zeile, he achieved the total synthesis of haemin, proving that the pigment consists of an iron atom surrounded by the porphyrin ring system.

\subsection*{Nobel Prize-winning Work}

The work for which Hans Fischer was awarded the Nobel Prize in Chemistry in 1930 was described by the Nobel Committee as: "for his researches into the constitution of haemin and chlorophyll and especially for his synthesis of haemin".

Fischer's investigations established the complete constitution of haemin, the red iron-containing pigment of blood, and showed that the green pigment chlorophyll is built from the same family of tetrapyrrole ring systems.
The total synthesis of haemin confirmed this structure beyond doubt.

\subsection*{Significance and Impact}

By establishing the chemistry of the porphyrins, Fischer created the foundation for understanding how oxygen is carried in the blood and how chlorophyll enables photosynthesis.
Porphyrin chemistry, developed in his school, remains central to organic and bioinorganic chemistry.

\subsection*{Later Career and Honors}

Fischer continued his research at the Technische Hochschule in Munich until the end of the Second World War.
He received the Liebig Medal in 1929, the Davy Medal of the Royal Society in 1937, and an honorary doctorate from Harvard University in 1936.
He died on 1945-03-31 in Munich, Germany.

\subsection*{Legacy}

The legacy of Hans Fischer endures through the lasting impact of his scientific contributions.
The chemistry of the porphyrins and bile pigments that he created continues to be built upon by researchers across the chemical and biological sciences.

\subsection*{Selected Publications}

\begin{itemize}

\item Fischer, H., \& Zeile, K. (1929). "Synthese des Hämatoporphyrins, Protoporphyrins und Hämins." Justus Liebigs Annalen der Chemie, 468, 98--116.

\item Fischer, H., \& Orth, H. (1934--1940). "Die Chemie des Pyrrols." Leipzig.

\end{itemize}

\clearpage

\section*{Chapter Summary}

The 1930 prize recognized Hans Fischer for his researches into the constitution of haemin and chlorophyll, especially his synthesis of haemin. His work revealed the porphyrin ring structure common to haemoglobin and chlorophyll.
